\documentclass{article}

% Umlaut-Support
\usepackage[ngerman]{babel}
\usepackage[utf8]{inputenc}
\usepackage[T1]{fontenc}

% AMS
\usepackage{amsmath}
\usepackage{amssymb}
\usepackage{amsthm}

% Koordinatensysteme
\usepackage{tkz-euclide}

% Grafiken und Bilder
\usepackage{graphicx}

% Seiteneinrichtung
\usepackage[a4paper, left=2cm, right=2cm, top=2cm, bottom=2cm]{geometry} % https://latex-kurs.blogspot.com/2012/09/latex-seitenrander-und-textwidth.html

% Befehle
\newcommand{\m}[1]{\begin{pmatrix}#1\end{pmatrix}}

\newcommand{\cspm}[1]{\begin{figure}[h]
    \begin{tikzpicture}
       \tkzInit[xmax=10,ymax=10,xmin=-10,ymin=-10]

        % Gitter
        \draw[step=5mm, gray, thin] (-5,-5) grid (5,5);

        % Koordinatensystem
        \foreach \x in {0,1,2,3,4,5}
            \draw (\x cm,1pt) -- (\x cm,-1pt) node[anchor=north] {$\x$};
        \foreach \y in {0,1,2,3,4,5}
            \draw (1pt,\y cm) -- (-1pt,\y cm) node[anchor=east] {$\y$};

        \foreach \x in {-1,-2,-3,-4,-5}
            \draw (\x cm,1pt) -- (\x cm,-1pt) node[anchor=north] {$\x$};
        \foreach \y in {-1,-2,-3,-4,-5}
            \draw (1pt,\y cm) -- (-1pt,\y cm) node[anchor=east] {$\y$};

        % Punkte und Vektoren
        #1
        
    \end{tikzpicture}
\end{figure}}

\newcommand{\csp}[1]{\begin{figure}[h]
    \begin{tikzpicture}
       \tkzInit[xmax=10,ymax=10,xmin=-10,ymin=-10]

        % Gitter
        \draw[step=5mm, gray, thin] (0,0) grid (10,10);

        % Koordinatensystem
        \foreach \x in {0,1,2,3,4,5,6,7,8,9,10}
            \draw (\x cm,1pt) -- (\x cm,-1pt) node[anchor=north] {$\x$};
        \foreach \y in {0,1,2,3,4,5,6,7,8,9,10}
            \draw (1pt,\y cm) -- (-1pt,\y cm) node[anchor=east] {$\y$};

        % Punkte und Vektoren
        #1
        
    \end{tikzpicture}
\end{figure}}

% Einstellungen
\setlength{\parindent}{0pt}

% Metadaten
\title{Verschiebung von Koordinatensystemen}
\author{Justus Seeck}
\date{\today}

% Dokumentenbeginn
\begin{document}
    \maketitle

    \tableofcontents

    \newpage

    \section{Einführung}

    \subsection{Vorwort}

    In der folgenden Arbeit wird lediglich die Verschiebung von Koordinatensystemen in der Ebene (also in zwei Dimensionen) behandelt.
    Die Methodik lässt sich jedoch auch auf die dritte Dimension übertragen. Auch die Anzahl der Verschiebungen ist beliebig, da man die
    hier verwendete Formel beliebig oft hintereinander anwenden kann, wird hier jedoch auf eine Verschiebung beschränkt.
    Die Verschiebung von Koordinatensystemen findet beispielweise in der Robotik Anwendung: Am ende eines Roboterarms befindet sich meist ein weiteres Koordinatensystem, "Toolsystem", welches sich am Ende des Roboterarms befindet. Dadurch gibt es in diesem System keine verschiebung wenn sich der Arm bewegt, sondern lediglich wenn sich das tool, oder dessen Ladnung bewegt.

    \subsection{Schreibweisen und Definitionen}
    
    \paragraph{Punkte}

    Ein Punkt $P$ wurde bisher in der Form $P(x,y)$ dargestellt. Dies wird in der folgenden Arbeit nicht mehr verwendet. Stattdessen wird
    der Punkt $P$ in der Form $\m{x \\ y}$ dargestellt. Dies ist eine sogenannte Matrix mit zwei Zeilen und einer Spalte.
    Die erste Zeile enthält die x-Koordinate, die zweite Zeile die y-Koordinate.
    Da im folgenden mit zwei verschiedenen Koordinatensystemen, dem Weltsystem $U$ und dem Toolsystem $T$, gearbeitet wird,
    werden Punkte im Bezug auf ein Koordinatensystem mit Index angegeben. Der Punkt $P$ im Bezug zum Weltsystem $U$ wird
    folgendermaßen dargestellt: ${[P]}_{U}$.
    Im Bezug zum Toolsystem lautet die Darstellung ${[P]}_{T}$.
    
    \section{Vektoren}

    Geometrische Vektoren sind Pfeile (oder Pfeilklassen). $\vec{v}$ ist ein Vektor.
    Vektoren werden durch ihre Länge in $x$-Richtung und $y$-Richtung definiert
    (bei drei Dimensionen zusätzlich in $z$-Richtung) und mit einem Pfeil über dem Buchstaben gekennzeichnet.
    Vektoren haben keine Feste Postition im Koordinatensystem. Sie können beliebig verschoben werden.
    Solange der Vektor alle seine Eigenschaften (Länge in $x$-, $y$- und ggf. $z$-Richtung) behält,
    ist er der gleiche Vektor. \textbf{Anhang: Vektoren 1}

    \subsection{Rechnen mit Vektoren}
    Vektoren können durch mathematische Operationen verändert werden. Multipliziert man beispielweise den Vektor $\m{2 \\ 0}$ mit $-1$,
    erhält man den Vektor $\m{-2 \\ 0}$. Man kan auch zwei Vektoren addieren. Aus $\m{2 \\ 3}$ und $\m{0 \\ 2}$ erhält man 

    % \subsection{Verschiebung von Punkten mittels Vektoren}
    % Um einen Punkt $P$ um einen Vektor $\vec{v}$ zu verschieben, wird der Punkt $P$ mit dem Vektor $\vec{v}$ multipliziert.
    

    % Anhänge

    \newpage

    \section{Anhänge}

    \subsection{Vektoren 1}
    \cspm{
        % Vektor 1
        \draw[->](-4,0.5) -- (-2,1.5) node[above] {$\vec{v}$};

        % Vektor 2
        \draw[->](2,2) -- (4,3) node[above] {$\vec{v}$};

        % Vektor 3
        \draw[->](-1,-3) -- (1,-2) node[above] {$\vec{v}$};
    }
    Koordinatensystem mit dem Vektor $\vec{v} = \m{2 \\ 1}$ in drei verschiedenen Positionen.

    \subsection{Vektoren 2}

    \begin{figure}[h]
        \centering
        \scalebox{0.5}{
            \includegraphics{./images/Mathemann.png}
        }
    \end{figure}

\end{document}


% Notizen und Hilfen

% Der Satz des Pytagoras lautet: $a^2 + b^2 = c^2$.
% Die PQ-Formel latuet:
    
% \[
%     x = - \frac{p}{2} \pm \sqrt{ \left ( \frac{p}{2}^2 \right ) - q }
% \]

% \[
%     \begin{pmatrix}
%         1 & 2 & 3 \\
%         4 & 5 & 6 \\
%         7 & 8 & 9

%     \end{pmatrix}
% \]
